\documentclass[11pt]{article}
\usepackage[letterpaper,margin=1in]{geometry}
\usepackage[T1]{fontenc}
\usepackage[utf8]{inputenc}
\usepackage{lmodern,microtype}
\usepackage{amsmath,amssymb,mathtools,bm,slashed}
\usepackage{xcolor,enumitem,booktabs,array,needspace}
\usepackage[most]{tcolorbox}
\usepackage{fancyhdr,titlesec,parskip,hyperref}
\definecolor{qftblue}{HTML}{245A8D}
\definecolor{qftgreen}{HTML}{2F6B4F}
\definecolor{qftgray}{HTML}{F4F6F8}
\hypersetup{colorlinks=true,linkcolor=qftblue,urlcolor=qftblue,pdfauthor={Nima Lashkari}}
\pagestyle{fancy}\fancyhf{}
\lhead{PHYS 662: Quantum Field Theory I}\cfoot{\thepage}
\setlength{\headheight}{14pt}
\titleformat{\section}{\Large\bfseries\color{qftblue}}{\thesection.}{0.6em}{}
\titleformat{\subsection}{\large\bfseries\color{qftgreen}}{\thesubsection}{0.6em}{}
\setlist[enumerate]{leftmargin=2em,itemsep=0.25em,topsep=0.35em}
\setlist[itemize]{leftmargin=2em,itemsep=0.2em,topsep=0.35em}
\newcommand{\dd}{\mathrm d}
\newcommand{\ii}{\mathrm i}
\newcommand{\eps}{\varepsilon}
\newcommand{\tr}{\operatorname{tr}}
\newcommand{\Tr}{\operatorname{Tr}}
\newcommand{\diag}{\operatorname{diag}}
\newcommand{\Lag}{\mathcal L}
\newcommand{\D}{\mathcal D}
\newcommand{\vac}{\lvert0\rangle}
\newcommand{\vbra}{\langle0\rvert}
\newcommand{\Tord}{\mathrm T}
\newcommand{\ket}[1]{\lvert#1\rangle}
\newcommand{\bra}[1]{\langle#1\rvert}
\newtcolorbox{keyidea}[1][]{enhanced,breakable,colback=qftgray,colframe=qftblue,
 boxrule=0.7pt,arc=1.5mm,title={#1},fonttitle=\bfseries,coltitle=white}
\newtcolorbox{remarkbox}[1][]{enhanced,breakable,colback=white,colframe=qftgreen,
 boxrule=0.6pt,arc=1.5mm,title={#1},fonttitle=\bfseries,coltitle=white}
\newcommand{\module}[1]{\clearpage\section{#1}}
\newcommand{\lecturefront}[4]{%
\rhead{Lecture #1}\hypersetup{pdftitle={PHYS 662 Lecture #1: #2}}
\begin{document}\pagenumbering{gobble}
\begin{titlepage}\centering\vspace*{1.05in}
{\scshape\Large PHYS 662 --- Quantum Field Theory I\par}
\vspace{0.35in}{\Huge\bfseries\color{qftblue}Lecture #1\par}
\vspace{0.2in}{\LARGE\bfseries #2\par}
\vspace{0.2in}{\large\color{qftgreen}#3\par}
\vspace{0.6in}{\large Nima Lashkari\par}{\large Fall 2026\par}
\vfill\begin{minipage}{0.86\textwidth}\small #4\par\medskip
Course page: \url{https://nimalashkari.com/quantum-field-theory-2026-part-i/}
\end{minipage}\vspace{0.5in}\end{titlepage}
\pagenumbering{roman}\setcounter{tocdepth}{1}\tableofcontents
\vfill\begin{remarkbox}[Place in the course]
Lectures 17--18 develop the quantum-mechanical construction. Lecture 19
extends Gaussian integration and source methods to scalar fields.
Lecture 20 introduces Grassmann integration and fermionic fields.
Problem Set 6 brings these four steps together.
\end{remarkbox}\clearpage\pagenumbering{arabic}}
\lecturefront{17}{Path Integrals in Quantum Mechanics I}
{From unitary evolution to a sum over histories}
{We derive the path integral from the operator formulation, fix its measure,
and recover the Schr\"odinger equation. Insertions and contact terms explain
how a description using ordinary paths remembers quantum operator ordering.}

\section{The evolution kernel and composition}
Lectures 15--16 described quantum fields by operators. Before changing
languages in field theory, we make the change for a single particle.
The path integral will represent the same evolution operator; it is not
an additional quantization postulate. In this lecture we keep \(\hbar\)
explicit and begin with one Cartesian coordinate.

For a time-independent, self-adjoint Hamiltonian, define
\begin{equation}
 K(q_f,t_f;q_i,t_i)=\bra{q_f}e^{-\ii H(t_f-t_i)/\hbar}\ket{q_i},
 \qquad \langle q|q'\rangle=\delta(q-q').                 \label{eq:kernel}
\end{equation}
If \(\psi_i(q_i)\) is the initial wavefunction, its later value is
\begin{equation}
 \psi_f(q_f)=\int\dd q_i\,K(q_f,t_f;q_i,t_i)\psi_i(q_i).
\end{equation}
Thus \(K\) is an amplitude density, of dimension inverse length. Its
absolute square at sharply specified endpoints is not, by itself, a
normalized transition probability. Probabilities are obtained from
normalized initial states and final measurements.

Insert the identity at an intermediate time \(t\). The group property
of evolution gives the exact composition law
\begin{equation}
 K(q_f,t_f;q_i,t_i)=\int\dd q\,
 K(q_f,t_f;q,t)K(q,t;q_i,t_i).                       \label{eq:composition}
\end{equation}
Iterating this identity expresses a finite-time amplitude in terms of
many short-time amplitudes. The sum over intermediate positions is the
origin of the sum over paths. No measurement of these intermediate
positions is performed: the alternatives are added as amplitudes.

\begin{keyidea}[Initial data and unitarity]
At coincident times, \(K(q_f,t_i;q_i,t_i)=\delta(q_f-q_i)\). Also,
\[
 \int\dd q_f\,K(q_f,t_f;q_i,t_i)^*K(q_f,t_f;q_j,t_i)
 =\delta(q_i-q_j).
\]
Any proposed path-integral normalization must satisfy these identities.
\end{keyidea}

For a time-dependent Hamiltonian, replace the exponential by the
time-ordered evolution operator. Composition still holds. We assume a
Hamiltonian bounded below when using the imaginary-time continuation
in Lecture 18; the real-time group law itself requires self-adjointness.

\module{Time slicing and the Hamiltonian symbol}
Set \(T=t_f-t_i>0\), \(\epsilon=T/N\), and
\(q_0=q_i,\ q_N=q_f\). Use the Fourier normalization
\begin{equation}
 \langle q|p\rangle=e^{\ii pq/\hbar},\qquad
 \langle p|p'\rangle=2\pi\hbar\,\delta(p-p'),\qquad
 1=\int\frac{\dd p}{2\pi\hbar}\ket p\bra p.
\end{equation}
For \(H=\hat p^2/(2m)+V(\hat q)\), the Trotter product gives
\begin{align}
 e^{-\ii H\epsilon/\hbar}
 &=e^{-\ii V(\hat q)\epsilon/\hbar}
   e^{-\ii\hat p^2\epsilon/(2m\hbar)}+O(\epsilon^2),\\
 \langle q_{j+1}|e^{-\ii H\epsilon/\hbar}|q_j\rangle
 &=\int\frac{\dd p_j}{2\pi\hbar}
 \exp\!\left\{\frac{\ii}{\hbar}
 [p_j(q_{j+1}-q_j)-\epsilon H(q_{j+1},p_j)]\right\}
 +O(\epsilon^2).
\end{align}
The limit is understood in the operator sense appropriate to the
Hamiltonian, or with a convergence regulator. Combining the short steps,
\begin{equation}
 K=\lim_{N\to\infty}\int\prod_{j=1}^{N-1}\dd q_j
       \prod_{j=0}^{N-1}\frac{\dd p_j}{2\pi\hbar}\,
 e^{\frac{\ii}{\hbar}\sum_{j=0}^{N-1}
 [p_j(q_{j+1}-q_j)-\epsilon H(q_{j+1},p_j)]}.       \label{eq:phase}
\end{equation}
There are \(N\) momentum integrals and \(N-1\) position integrals.
The endpoints are fixed, so integrating them would calculate a different
matrix element. The compact notation is
\(K=\int\D p\,\D q\,e^{\ii\int(p\dot q-H)\dd t/\hbar}\).

\subsection{Why ordering survives the continuum limit}
For a general Hamiltonian containing mixed products, the ordinary
function in \eqref{eq:phase} is an \emph{operator symbol}. Putting all
\(\hat q\)'s to the left of all \(\hat p\)'s gives the displayed endpoint
prescription. Weyl ordering instead uses a midpoint symbol.
For example,
\begin{equation}
 \frac12(\hat q\hat p+\hat p\hat q)
 =\hat q\hat p-\frac{\ii\hbar}{2}.
\end{equation}
Its \(qp\)-ordered symbol is \(qp-\ii\hbar/2\), whereas its Weyl
symbol is \(qp\). The difference per time slice is proportional to
\(\epsilon\); after \(N\) slices it need not disappear.
\begin{remarkbox}[A consistent prescription]
Choose the operator ordering, its symbol, and the time discretization
together. For \(p^2/(2m)+V(q)\) the usual endpoint and midpoint
constructions give the same continuum Hamiltonian under the standard
regularity assumptions. This special case does not eliminate ordering
issues for arbitrary mixed Hamiltonians.
\end{remarkbox}

\module{Integrating momentum: the configuration measure}
Complete the square in each momentum integral:
\begin{align}
 p\Delta q-\epsilon\frac{p^2}{2m}
 &=-\frac{\epsilon}{2m}\left(p-\frac{m\Delta q}{\epsilon}\right)^2
   +\frac{m(\Delta q)^2}{2\epsilon},\\
 \int\frac{\dd p}{2\pi\hbar}
 e^{\frac{\ii}{\hbar}(p\Delta q-\epsilon p^2/2m)}
 &=\left(\frac{m}{2\pi\ii\hbar\epsilon}\right)^{1/2}
 e^{\frac{\ii m(\Delta q)^2}{2\hbar\epsilon}}.       \label{eq:gaussian}
\end{align}
The oscillatory Gaussian is defined by continuation from
\(\epsilon\to\epsilon-\ii0\). This choice also fixes the square-root
branch. It is not a positive real Gaussian before continuation.

Consequently,
\begin{equation}
 K=\lim_{N\to\infty}\left(\frac{m}{2\pi\ii\hbar\epsilon}\right)^{N/2}
 \int\prod_{j=1}^{N-1}\dd q_j\,
 e^{\frac{\ii}{\hbar}\sum_j
 [\frac{m(q_{j+1}-q_j)^2}{2\epsilon}-\epsilon V(q_{j+1})]}.
                                                        \label{eq:config}
\end{equation}
We abbreviate this regulated limit as
\begin{equation}
 K(q_f,t_f;q_i,t_i)=\int_{q_i}^{q_f}\D q\,
 e^{\ii S[q]/\hbar},\qquad
 S[q]=\int_{t_i}^{t_f}\dd t\left(\frac m2\dot q^2-V(q)\right).
\end{equation}
The symbol \(\D q\) includes the prefactor in \eqref{eq:config}.
Dropping it loses the short-time delta function and the composition law.
For \(d\) independent Cartesian coordinates its power becomes \(Nd/2\).

\subsection{What is a path in this expression?}
The finite-dimensional integral, together with its limiting prescription,
defines the notation. It does not assume that every contributing path
has an ordinary derivative. After continuation to Euclidean time, a
typical short increment has size \(\Delta q\sim\sqrt{\hbar\epsilon/m}\),
so \(\Delta q/\epsilon\) grows as the slicing is refined. Smooth paths
are useful saddle points, not representative integration histories.

\begin{keyidea}[The measure is part of the physics]
An action written without a measure and a time-slicing prescription is
incomplete. In particular, transforming to curvilinear coordinates
requires transforming the measure and the regulated kinetic operator.
Simply replacing \(\dot q^2\) by a coordinate-dependent metric can miss
ordering terms.
\end{keyidea}
The continuum limit in this lecture is controlled by a specified quantum
Hamiltonian. In field theory an additional spatial regulator is needed;
removing it while holding physical parameters fixed is a separate issue.

\module{An exact test: the free particle}
For \(V=0\), separate the straight path and its endpoint-fixed fluctuations:
\begin{equation}
 q_j=q_i+\frac jN(q_f-q_i)+\xi_j,\qquad\xi_0=\xi_N=0.
\end{equation}
The cross term telescopes, leaving
\begin{equation}
 \sum_{j=0}^{N-1}(q_{j+1}-q_j)^2
 =\frac{(q_f-q_i)^2}{N}+\bm\xi^{\mathsf T}A_{N-1}\bm\xi,
 \quad
 A_n=\begin{pmatrix}
 2&-1&0&\cdots\\-1&2&-1&\cdots\\0&-1&2&\ddots\\
 \vdots&\vdots&\ddots&\ddots
 \end{pmatrix}.
\end{equation}
Expanding the determinant along its first row gives
\(d_n=2d_{n-1}-d_{n-2}\), with \(d_0=1,d_1=2\).
Thus \(\det A_{N-1}=N\). The \(N-1\) Gaussian integrations contribute
\((2\pi\ii\hbar\epsilon/m)^{(N-1)/2}/\sqrt N\).
Multiplying by the measure in \eqref{eq:config} gives
\begin{keyidea}[Free kernel in \(d\) spatial dimensions]
\begin{equation}
 K_0(\bm q_f,\bm q_i;T)=
 \left(\frac{m}{2\pi\ii\hbar T}\right)^{d/2}
 \exp\!\left[\frac{\ii m|\bm q_f-\bm q_i|^2}{2\hbar T}\right],
 \qquad T>0.                                         \label{eq:free}
\end{equation}
The exponent is the classical action. The prefactor is the exact
contribution of fluctuations about that path.
\end{keyidea}

An independent derivation inserts a momentum identity:
\begin{equation}
 K_0=\int\frac{\dd^d p}{(2\pi\hbar)^d}
 e^{\frac{\ii}{\hbar}[\bm p\cdot(\bm q_f-\bm q_i)-\bm p^2T/(2m)]}.
                                                        \label{eq:freefourier}
\end{equation}
This form immediately proves the delta-function initial condition.
It also proves composition: integrating an intermediate position sets
the two momenta equal, and their evolution phases multiply.

For fixed \(T\), \(K_0\) is nonzero at arbitrarily large spatial
separation. This is nonrelativistic quantum mechanics, with no finite
limiting speed built into \(H_0\). Later, relativistic field locality
will be a statement about commutators. It does not require the Feynman
two-point function to vanish outside the light cone.

\module{The Schr\"odinger equation and its Green function}
The operator definition already implies
\begin{equation}
 (\ii\hbar\partial_{t_f}-H_{q_f})K(q_f,t_f;q_i,t_i)=0,
 \qquad H_q=-\frac{\hbar^2}{2m}\partial_q^2+V(q).       \label{eq:schrod}
\end{equation}
The path integral supplies a useful local check. Set
\(q_i=q-\xi\) over a short interval. Acting on a smooth test
wavefunction, the regulated Gaussian moments are
\begin{equation}
 \langle1\rangle_\epsilon=1,\qquad
 \langle\xi\rangle_\epsilon=0,\qquad
 \langle\xi^2\rangle_\epsilon=\frac{\ii\hbar\epsilon}{m}.
\end{equation}
Expanding the potential phase and the wavefunction therefore gives
\begin{align}
 \psi(q,t+\epsilon)
 &=\int\dd\xi\,K(q,t+\epsilon;q-\xi,t)\psi(q-\xi,t)\\
 &=\psi(q,t)+\frac{\ii\hbar\epsilon}{2m}\partial_q^2\psi(q,t)
   -\frac{\ii\epsilon}{\hbar}V(q)\psi(q,t)+O(\epsilon^2).
\end{align}
Multiplying by \(\ii\hbar/\epsilon\) recovers the usual equation.
The Gaussian second moment is essential: setting every short path equal
to a classical path would miss the kinetic differential operator.

\subsection{Where the delta-function source belongs}
Evolution is defined for either sign of \(T\). To construct a retarded
inverse of the differential operator, introduce instead
\begin{equation}
 K_R=\theta(T)K,\qquad G_R=\frac{1}{\ii\hbar}\theta(T)K.
\end{equation}
Differentiating the step function yields
\begin{align}
 (\ii\hbar\partial_{t_f}-H_{q_f})K_R
 &=\ii\hbar\,\delta(T)\delta(q_f-q_i),\\
 (\ii\hbar\partial_{t_f}-H_{q_f})G_R
 &=\delta(T)\delta(q_f-q_i).                          \label{eq:green}
\end{align}
These formulas have different normalizations because they define
different objects. A delta-function initial condition for \(K\) does
not mean that its homogeneous evolution equation has a source.

\begin{remarkbox}[Notation to carry forward]
\(K\) denotes the evolution kernel; \(G_R\) denotes its unit-source
retarded Green function. In field theory, \(D_F\) will retain the
Lecture 16 convention \(D_F=\langle0|\Tord\phi\phi|0\rangle\),
including its factor of \(\ii\) in momentum space.
\end{remarkbox}

\module{Insertions and time ordering}
Choose \(t_i<t_1<t_2<t_f\). Inserting position resolutions on the
two specified slices gives
\begin{align}
 &\int_{q_i}^{q_f}\D q\,q(t_2)q(t_1)e^{\ii S[q]/\hbar}\nonumber\\
 &\quad=\bra{q_f}U(t_f,t_2)\hat q\,U(t_2,t_1)\hat q\,
 U(t_1,t_i)\ket{q_i}.                              \label{eq:insert}
\end{align}
Taking the Heisenberg reference time to be \(t_i\), this equals
\(\bra{q_f}U(t_f,t_i)\hat q_H(t_2)\hat q_H(t_1)\ket{q_i}\).
For either ordering of the times, the path integral represents the
\emph{time-ordered} operator product. The later operator is nearer
the final bra. This follows from composition, rather than from a
new rule imposed on the integral.

The path values are ordinary commuting numbers, but the quantum
operators do not commute. The apparent tension disappears because
the map from histories to operators already includes time ordering.
An unordered or out-of-time-order correlator generally needs a
different arrangement of evolution segments.

\subsection{A source collects all insertions}
Introduce an ordinary function \(J(t)\) supported away from the endpoints:
\begin{equation}
 K[J]=\int_{q_i}^{q_f}\D q\,
 \exp\!\left\{\frac{\ii}{\hbar}\left[S[q]+\int\dd t\,J(t)q(t)\right]\right\}.
\end{equation}
Then
\begin{equation}
 \left.\left(\frac{\hbar}{\ii}\right)^n
 \frac{\delta^n K[J]}{\delta J(t_1)\cdots\delta J(t_n)}\right|_{J=0}
 =\int_{q_i}^{q_f}\D q\,q(t_1)\cdots q(t_n)e^{\ii S/\hbar}.
\end{equation}
These are endpoint matrix elements. They become vacuum expectation
values only after a vacuum projection or an appropriate vacuum boundary
wavefunction is supplied. Lecture 18 explains this step.

Momentum insertions can be handled directly in the phase-space
integral. Composite insertions such as \(pq\), or fields at exactly
coincident times, retain the ordering prescription from the time
slices. Replacing \(p\) by \(m\dot q\) inside every correlation function
without tracking contact terms is not generally justified.

\begin{keyidea}[A useful distinction]
Endpoint conditions select the states. The arrangement of insertions
selects the operator ordering. The action alone specifies neither.
\end{keyidea}

\module{Schwinger--Dyson identities and contact terms}
On a finite lattice, the integral of a total derivative vanishes when
the regulated integrand has suitable boundary behavior. Vary an
interior path value while keeping both endpoints fixed. In continuum
notation the resulting identity is
\begin{align}
 0&=\int\D q\,\frac{\delta}{\delta q(t)}
       \left(F[q]e^{\ii S[q]/\hbar}\right),\\
 \int\D q\,F[q]\frac{\delta S}{\delta q(t)}e^{\ii S/\hbar}
 &=\ii\hbar\int\D q\,\frac{\delta F}{\delta q(t)}e^{\ii S/\hbar}.
                                                        \label{eq:sd}
\end{align}
The sign is \(+\ii\hbar\), since \(-\hbar/\ii=\ii\hbar\).
For our action,
\begin{equation}
 \frac{\delta S}{\delta q(t)}=-m\ddot q(t)-V'(q(t)).
\end{equation}
With \(F=1\), the classical equation holds in the corresponding
matrix element. With \(F=q(t')\), however,
\begin{equation}
 \int\D q\,[m\ddot q(t)+V'(q(t))]q(t')e^{\ii S/\hbar}
 =-\ii\hbar\delta(t-t')K[0].                         \label{eq:contact}
\end{equation}
One must keep the contact term when differentiating a time-ordered
correlator. Away from \(t=t'\), it reduces to the equation of motion.

\subsection{Operator check of the contact sign}
For the harmonic oscillator, let
\(C(t,t')=\langle0|\Tord q(t)q(t')|0\rangle\).
The first derivative of time ordering has no contact term because
\([q,q]=0\). The second derivative produces
\([\dot q(t),q(t)]=-[q,p]/m=-\ii\hbar/m\). Hence
\begin{equation}
 m(\partial_t^2+\omega^2)C(t,t')=-\ii\hbar\delta(t-t').
\end{equation}
Indeed \(C=\hbar e^{-\ii\omega|t-t'|}/(2m\omega)\) has precisely
the required derivative jump. Thus integration by parts in the path
integral remembers the canonical commutator.

For several insertions, \(\delta F/\delta q(t)\) is a sum of delta
functions, one for every insertion containing the varied coordinate.
This hierarchy is the Schwinger--Dyson hierarchy. It organizes exact
relations among correlation functions even when the integral cannot
be evaluated explicitly. In interacting theories it does not usually
close on a finite number of correlators.

\module{Stationary phase and the classical limit}
The phase changes rapidly when the action is large compared with
\(\hbar\). Around a smooth stationary path write \(q=q_{\rm cl}+\xi\),
with \(\xi(t_i)=\xi(t_f)=0\). The first variation is
\begin{equation}
 \delta S=[m\dot q\,\xi]_{t_i}^{t_f}
 -\int\dd t\,[m\ddot q+V'(q)]\xi.
\end{equation}
It vanishes for a classical solution satisfying the endpoint conditions.
The next terms are
\begin{equation}
 S[q]=S[q_{\rm cl}]+\frac12\int\dd t\,
 \xi\left[-m\frac{\dd^2}{\dd t^2}-V''(q_{\rm cl})\right]\xi
 +O(\xi^3).                                           \label{eq:saddle}
\end{equation}
The quadratic functional integral supplies a fluctuation determinant.
For the free particle it is exact; for a general potential it is the
leading semiclassical correction. If several classical paths connect
the endpoints, their amplitudes must be summed, including their
relative phases. Near a zero eigenvalue of the fluctuation operator,
the simplest Gaussian approximation fails and needs refinement.

\begin{remarkbox}[Stationary action, not always minimum action]
Real-time classical paths are stationary points of the action.
The quadratic form may have negative directions, and destructive
interference is the relevant mechanism. A positive Euclidean weight,
introduced next lecture, supports a different saddle-point intuition.
\end{remarkbox}

The derivation explains the classical equations without claiming that
all paths are classical or that the quantum measure becomes an ordinary
uniform measure on smooth trajectories. It also explains why a
semiclassical approximation and an expansion in a weak interaction are
different approximations: the first organizes powers of \(\hbar\),
whereas the second organizes powers of a coupling.

\subsection*{Questions to check understanding}
\begin{enumerate}
\item Why does a discretized kernel have one more momentum integral
than position integral?
\item Which factor in the free kernel would be missed by evaluating
only the classical action?
\item How can two Hamiltonian symbols differing by \(O(\hbar)\)
describe the same operator under different ordering prescriptions?
\item What distinguishes \(K\), \(\theta(T)K\), and \(G_R\)?
\item Why do commuting path values generate time-ordered quantum operators?
\item Derive the contact term in \eqref{eq:contact} in both languages.
\end{enumerate}
Next lecture uses the same construction for perturbations, scattering,
vacuum projection, and imaginary-time evolution.
\end{document}
